\begin{itemize}
\item \textbf{NORMAL.} \emph{Numérica:} En vibración libre RK4 reproduce la solución cerrada con error máximo 6.161$\times10^{-7}$ m y Heun con 0.00123 m; los órdenes observados se acercan a 4 y 2. Simpson y Romberg integran x$^{2}$ con muchas menos evaluaciones que el trapecio para el mismo error. \emph{Problema:} Sin viento, la oscilación inicial de 0.1 m decae con envolvente e\^{}($-\zeta\omega_{0}$t): a los 20 s queda 0.0135 m. El amortiguamiento controla cuánto dura el balanceo.
\item \textbf{ESTRÉS.} \emph{Numérica:} Con h = 0.5 s ($\omega_{0}$h = 1) Heun tiene |R(h$\lambda$)| = 1.051 > 1: amplifica la oscilación en cada paso y su error máximo llega a 295 m, mientras RK4 (|R| = 0.9467) se queda en 0.286 m. Al reducir h, Heun vuelve a ser estable y ambos recuperan sus órdenes 2 y 4. La diferencia Heun/RK4 es de estabilidad, no sólo de precisión. \emph{Problema:} La ráfaga con $\Omega$/$\omega_{0}$ = 0.95 produce una amplitud permanente de 2.296 m, 7.346 veces la deflexión estática F$_{0}$/k. Cerca de la resonancia la amplitud depende sobre todo del amortiguamiento y de la frecuencia de la ráfaga, no de la masa por sí sola. Amplitudes de ese orden superan el rango de validez de un modelo lineal de pequeñas oscilaciones.
\item \textbf{MAL COND..} \emph{Numérica:} Con $\zeta$ = 0.005 el problema es sensible: $\kappa$(X; $\Omega$) = 99, de modo que un error relativo de 0.1 \% en $\Omega$ se convierte en $\approx$ 9.9 \% en la amplitud. Numéricamente, el error de fase de Heun se acumula durante 300 s (error máx. 0.186 m) mientras RK4 queda en 1.026$\times10^{-4}$ m; ambos convergen al refinar h. La sensibilidad es del modelo, no un fallo de los métodos. \emph{Problema:} La ráfaga con $\Omega$/$\omega_{0}$ = 0.995 produce una amplitud permanente de 0.7393 m, 70.98 veces la deflexión estática F$_{0}$/k. Cerca de la resonancia la amplitud depende sobre todo del amortiguamiento y de la frecuencia de la ráfaga, no de la masa por sí sola. Amplitudes de ese orden superan el rango de validez de un modelo lineal de pequeñas oscilaciones.
\end{itemize}
